%! TeX program = lualatex
% \documentclass[twoside, a4paper, 12pt]{book}
\documentclass[a4paper, 12pt]{article}

\title{Hiperbole tańczą, tańczą hiperbole}
\author{\color{subtext1}Weronika Jakimowicz}
\date{2026}

\usepackage[utf8]{inputenc}
\usepackage[T1]{fontenc}

\usepackage{fontspec}

\usepackage{float}

\usepackage[default]{cantarell}


\usepackage[
  a4paper,
  total={170mm, 247mm},
  left=20mm,
  top=25mm,
  % showframe
]{geometry}

\usepackage[skip=3mm]{parskip}
\renewcommand{\baselinestretch}{1.2}
\selectfont


\usepackage{amsfonts}
\usepackage{amsmath}
\usepackage{amssymb}
\usepackage{amsthm}

\usepackage{tikz}
\usepackage{tikz-cd}
\usetikzlibrary{decorations.pathreplacing,calc,math,backgrounds, arrows}
\usetikzlibrary{positioning, decorations.text, decorations.pathmorphing}

\usepackage[nolessnomore,noequal,noparenthesis,nospecials,noplusnominus,noexclam,italic]{mathastext}

\DeclareMathOperator{\Spec}{Spec}
\DeclareMathOperator{\rank}{rank}
\DeclareMathOperator{\diam}{diam}
\DeclareMathOperator{\vol}{Vol}
\DeclareMathOperator{\supp}{supp}
\DeclareMathOperator{\Func}{Func}
\DeclareMathOperator{\dom}{dom}
\DeclareMathOperator{\ord}{ord}
\DeclareMathOperator{\Div}{Div}
\DeclareMathOperator{\Ob}{Ob}
\DeclareMathOperator{\Hom}{Hom}
\DeclareMathOperator{\Nat}{Nat}

\DeclareMathOperator{\Aut}{Aut}

\DeclareMathOperator{\Image}{Im}
\renewcommand{\Im}{\Image}
\DeclareMathOperator{\Real}{Re}
\renewcommand{\Re}{\Real}

\DeclareMathOperator{\alg}{\mathfrak{alg}}

\newcommand{\Z}{\mathbb{Z}}
\newcommand{\N}{\mathbb{N}}
\newcommand{\R}{\mathbb{R}}
\newcommand{\Q}{\mathbb{Q}}
\newcommand{\C}{\mathbb{C}}
\newcommand{\A}{\mathbb{A}}
\newcommand{\Proj}{\mathbb{P}}
\newcommand{\Hyp}{\mathbb{H}}

\newcommand{\Cc}{\mathcal{C}}
\newcommand{\Dd}{\mathcal{D}}
\newcommand{\Oo}{\mathcal{O}}

\newcommand{\mm}{\mathfrak{m}}

\let\epsilon\varepsilon

\renewcommand{\textbf}[1]{{\bfseries\boldmath#1}}

\usepackage{xcolor}

\definecolor{rosewater}{HTML}{dc8a78}
\definecolor{flamingo}{HTML}{dd7878}
\definecolor{pink}{HTML}{ea76cb}
\definecolor{mauve}{HTML}{8839ef}
\definecolor{red}{HTML}{d20f39}
\definecolor{maroon}{HTML}{e64553}
\definecolor{peach}{HTML}{fe640b}
\definecolor{yellow}{HTML}{df8e1d}
\definecolor{green}{HTML}{40a02b}
\definecolor{teal}{HTML}{179299}
\definecolor{sky}{HTML}{04a5e5}
\definecolor{sapphire}{HTML}{209fb5}
\definecolor{blue}{HTML}{1e66f5}
\definecolor{lavender}{HTML}{7287fd}
\definecolor{subtext2}{HTML}{4c4f69}
\definecolor{subtext1}{HTML}{5c5f77}
\definecolor{subtext0}{HTML}{6c6f85}
\definecolor{overlay2}{HTML}{7c7f93}
\definecolor{overlay1}{HTML}{8c8fa1}
\definecolor{overlay0}{HTML}{9ca0b0}

\definecolor{text}{HTML}{000000}
\color{text}

\newtheorem{lemma}{Lemat}
\newtheorem{theorem}{Twierdzenie}

\theoremstyle{definition}
\newtheorem{example}{Przykład}

\renewcommand{\figurename}{Rysunek}




\begin{document}
\maketitle
\thispagestyle{empty}

Badanie form kwadratowych -> ciekawe pytanie z teorii liczb. Nie umiem pisać wstępóóóów.

\section{Przestrzeń form kwadratowych}

Interesują nas formy kwadratowe na $\R^2$ z ustaloną bazą. Każda taka forma zapisuje się jako
$$Q(x,y)=ax^2+bxy+cy^2,$$
przypisujemy jej więc symetryczną macierz
$$A_Q=\begin{bmatrix}
  a & \frac{b}{2} \\ 
  \frac{b}{2} & c
\end{bmatrix}$$
Przestrzeń niezdegenerowanych form kwadratowych jest trójwymiarowa, ale możemy uprościć jej badanie przez pytanie o przestrzeń
$$ax^2+bxy+cy^2=0,$$
która jest niezmiennicza na skalowanie. Innymi słowy interesuje nas projektywizacja przestrzeni niezdegenerowanych form kwadratowych, czyli $\R P^2$.

Pierwszym niezmiennikiem formy kwadratowej $Q$ jest wyznacznik jej macierzy $A_Q$ z dokładnością do skalowania, $\det(A_Q)=ac-\frac{b^2}{4}$. Pojawia się pierwsza klasyfikacja rzeczywistych form kwadratowych:
\begin{enumerate}
  \item $\det(A_Q)>0$, wtedy dyskryminant równania $ax^2+bxy+cy^2=0$ jest ujemny, więc nie posiada ono rozwiązań rzeczywistych poza $(0,0)$. Dostajemy więc równanie postaci $(x-zy)(x-\overline{z}y)=0$ dla pewnej liczby zespolonej $z$. 
  \item $\det(A_Q)=0$, wtedy dyskryminant równania również jest zerowy więc jest jedno podwójne równanie. Dostajemy wtedy tak zwaną \emph{podwójną prostą} $(x-ay)^2=0$.
  \item $\det(A_Q)<0$, wtedy dyskryminant równania jest dodatni, widzimy więc dwie liniowo niezależne proste $(x-ay)(x-by)=0$.
\end{enumerate}

Geometrycznie formy o zerowym wyznaczniku definiują w $\R P^2$ krzywą
$$b^2-4ac=0$$
która jest okręgiem. Zmieniamy współrzędne na $b=2x$, $a=y+z$ oraz $c=z-y$ i otrzymujemy równanie okręgu o promieniu $z$
$$4x^2=4(z^2-y^2)\to x^2+y^2=z^2.$$
Jego wnętrze, topologicznie dysk, to formy spełniające $\det(A_Q)>0$. Możemy na tej przestrzeni określić odległość między dwoma punktami $A$, $B$ jako dwustosunek między punktami $P$ i $Q$ będącymi punktami przecięcia prostej przechodzącej przez $A$ i $B$ z krzywą $\det(A_Q)=0$. Ta konstrukcja to dokładnie model Kleine'a płaszczyzny hiperbolicznej.

\section{Klasyfikowanie form przy pomocy płaszczyzny hiperbolicznej}

Górna półpłaszczyzna z metryką $ds^2=\frac{dx^2+dy^2}{y^2}$ jest modelem płaszczyzny hiperbolicznej Poincarego. Wygodnie jest o niej myśleć jako o górnej półpłaszczyźnie zespolonej $H^2$ z grupą automorfizmów
$$\Aut(H^2)=\{z\mapsto \frac{az+b}{cz+d}\;:\;a,b,c,d\in \R\}\cong PSL(2, \R).$$


Niech $Q(x,y)=ax^2+bxy+cy^2$ będzie formą z obszaru ograniczonego przez okrąg $\det(A_Q)=0$. Wtedy $Q(x,y)=(x-zy)(x-\overline{z}y)$ dla pewnej liczby zespolonej $z$ takiej, że $\Im(z)>0$.

\begin{lemma}
  Dla formy kwadratowej $Q$ i liczby zespolonej $z_0$ jak wyżej istnieje dokładnie jedna liczba 
  $$z_1\in \{z\;:\;-\frac{1}{2}< \Re(z)\leq \frac{1}{2}\;\land\;|z|\geq1\}=P$$
  taka że istnieje element $M\in PSL(2, \Z)$ taki, że $z_0=Mz_1$.
\end{lemma}

\begin{proof}
  Niech $n=\lfloor\Re(z_0)\rfloor$. Rozważmy działanie macierzy 
  $$M_n=\begin{bmatrix}1&-n\\0&1\end{bmatrix}.$$
  Wtedy oczywiście $-\frac{1}{2}<\Re(M_n(z_0))\leq\frac{1}{2}$. Może się jednak przydarzyć, że $\Im(M_n(z_0))$ jest na tyle małe, że $|z_0|< 1$. Rozważamy wtedy macierz 
  $$A=\begin{bmatrix}0&-1\\1&0\end{bmatrix}$$
  wtedy punkt $z_0'=A(z_0)=\frac{-1}{z_0}$ leży poza okręgiem jednostkowym. Powtarzając te dwie operacje w skończenie wielu krokach trafimy do punktu w $P$.

  Dla jedyności wystarczy pokazać, że jeśli $z_1$, $z_2\in P$ to nie istnieje macierz $M\in PSL(2, \Z)-\{Id\}$ taka, że $z_1=Mz_2$. Załóżmy nie wprost, że istnieje
  $$M=\begin{bmatrix}a & b\\c&d\end{bmatrix}\in PSL(2, Z)$$
  takie, że $z_1=Mz_2$. Bez starty ogólności $\Im(z_1)\geq\Im(z_2)$. Wtedy 
  $$\Im(Mz_1)=\frac{\Im(z_1)}{|cz_1+d|^2}$$
  czyli $|cz_1+d|\leq 1$. Wtedy
  \begin{itemize}
    \item $c=0\implies d=1\implies a=1\implies Mz_1=z_1+b\implies b=0\implies M=Id\implies z_1=z_2$ co daje sprzeczność
    \item $c=\pm1\implies |z_1+d|\leq 1$. Jeśli $d=0$ to wtedy $|z_1|\leq 1$ i $z_1\not\in P$ co jest sprzeczne. Jeśli $d=\pm 1$ to $|z_1+d|^2=|z_1|+d^2+2x_1d>|z_1|^2+d^2-|d|>|z_1|^2>1$.
    \item $|c|\geq2\implies |cz_1+d|^2\geq |cz_1|^2=|c|^2|z_1|^2\geq|c|\geq 2$.
  \end{itemize}
  W takim razie $z_1=z_2$ i $M=Id$.
\end{proof}

W takim razie dobrze określonym niezmiennikiem formy $Q$ takiej, że $\det(A_Q)>0$ jest liczba zespolona $z\in P$. Powstaje pytanie czy istnieje analogiczny sposób klasyfikacji form o ujemnym wyznaczniku macierzy?

Dla form o ujemnym wyznaczniku, czyli lądujących poza okręgiem $\det(A_Q)=0$, możemy poprowadzić dwie proste styczne do tego okręgu przechodzące przez interesującą nas formę. To daje dwa punkty na okręgu, które możemy połączyć cięciwą w okręgu. Interpretuje się to jako geodezyjna na płaszczyźnie hiperbolicznej, czyli półokrąg oparty na osi rzeczywistej.

PRZYKŁAD

Rozważam $x^2+4xy+y^2$, czyli punkt $(1, 4, 1)$. Szukam prostych w $\R P^2$, dostaję
$$(1+t(2+\sqrt{3});4+t;1)$$
$$(1+t(2-\sqrt{3});4+t;1)$$
pierwsza daje formę $(y-(2+\sqrt{3})x)^2$, druga to $(y-(2-\sqrt{3})x)^2$

cyk mam okrąg
$$(x-2)^2+y^2=3\implies |z-2|=\sqrt{3}$$
i w sumie to utknęłam, bo dla $x=\frac{1}{2}$ idealnie przecinam okrąg $|z|=1$

czuję tutaj ułamki łańcuchowe? akurat dla $2+\sqrt{3}$ ląduję po inwersji w $\sqrt{3}-2$ i po jeszcze jednej wracam do $2+\sqrt{3}$ i nie robię po drodze odbijania o $x=\sqrt{1}{2}$


Jak mam $(1;3;1)$ to punkty na $OX$ to
dobra, mam
$(1+t(3\pm\sqrt{5});3+2t;1)$
trzeba rozwiązać dla jakich $t$ spełnia $b^2-4a=0$





% Formy kwadratowe na płaszczyźnie zapisane w pewnej ustalonej bazie tworzą przestrzeń trójwymiarową:
% $$ax^2+bxy+cy^2\mapsto (a,b,c),$$
% ale ponieważ interesują nas rozwiązania równania
% $$ax^2+bxy+cy^2=0$$
% to możemy przejść do projektywizacji przestrzeni form kwadratowych, czyli w przypadku form o współczynnikach rzeczywistych - przestrzeni $\R P^2$.
%
% Każdej takiej formie możemy przypisujemy macierz
% $$\begin{bmatrix}
%   a & \frac{b}{2} \\ 
%   \frac{b}{2} & c
% \end{bmatrix}$$
%
%
%
% \section{Model płaszczyzny hiperbolicznej}
%
% górna półpłaszczyzna zespolona $H=\{z\in\C\;:\;\Im(z)\geq0\}$
%
% mamy formę kwadratową
% $$ax^2+bxy+cy^2$$
% czyli macierz
% $$\begin{bmatrix}
%   a & \frac{b}{2} \\ 
%   \frac{b}{2} & c
% \end{bmatrix}$$
% z wielomianem charakterystycznym
% $$t^2+(a+c)t+(ac-\frac{b^2}{4})$$
% ta macierz ma wartości własne - pytanie ile rzeczywistych
%
% wypadałoby przypomnieć mapping class group i torusika?
%
% % to się wogóle nazywa conic section?
%
%
% elipsa o objętości $1$ i środku $0$:
% $$ax^2+bxy+cy^2=1$$
% tu nie jestem pewna co zrobić 
% no potrafię to zmienić na $x^2+y^2=0$ bo jak przerzucę $1$ to będę musiała przy reparametryzacji podoklejać $i$
%
%
% para prostych przechodzących przez $0$:
% $$0=(x-ay)(x-by)=x^2-(a+b)xy+aby^2$$
% to potrafię zmienić tak, żeby było
% $$0=(x'-y')(x'+y')$$
% na obrazku widać
%
%
% prosta podwójna:
% $$0=(x-ay)^2=x^2-2axy+a^2y^2$$
% tutaj faktycznie $(x,y)$ oraz $(-x, -y)$ daje to samo i wogóle skalowanie mi nic nie zmienia jak na $\R P^2$
%
%
%
% jeśli mam elipsę
%

\end{document}
